% ---------------------------------------------------------------------------
% Preamble for the PDF build. Loaded via `include-in-header` in _quarto.yml.
%
% This replaced a ~380 line tufte.tex on 4 Aug 2026. Almost all of that file
% was working around the `tufte-book` class rather than styling anything:
% patching \allcaps against the LaTeX kernel, re-defining \subsubsection,
% widening `Shaded` by \@tufte@overhang, aliasing `figure` to `figure*`.
%
% The decisive problem was that tufte-book defines \sidenote and \marginnote
% itself, which collides with the `sidenotes` and `marginnote` packages Quarto
% emits. The old file resolved that by declaring both packages already loaded
% so Quarto would skip them - which silently disabled `.column-margin`,
% `cap-location: margin` and `reference-location: margin`. The class was
% blocking the margin features it was chosen for.
%
% So: plain KOMA `scrbook`, wide margins from `geometry`, and Quarto's own
% margin machinery. What is left here is the two things that are genuinely
% ours.
% ---------------------------------------------------------------------------

% --- Image sizing ----------------------------------------------------------
%
% Images arrive three ways:
%
%   1. `width=100%` in the .qmd  -> \includegraphics[width=1\linewidth]
%   2. `width=70%`  in the .qmd  -> \includegraphics[width=0.7\linewidth]
%   3. no width at all           -> \pandocbounded{\includegraphics{...}},
%      which only ever *shrinks* an image to fit, never grows it. A wide
%      screenshot still fills the whole measure.
%
% THE .QMD WINS. Our width goes FIRST and the passed options `#1` come after,
% and graphicx takes the last value given for a key. So an explicit `width=`
% in the source overrides this, and only images that asked for nothing fall
% back to \rbpfigscale.
%
% This was tried the other way round first, forcing every image to the full
% measure, which is why the screenshots came out enormous. Dropping them all
% to one small fraction is not the fix either: at 0.5 the purpose-built
% teaching diagrams became too small to read the labels on, which is why they
% ask for `width=100%` in the .qmd.
%
% THE KNOB: \rbpfigscale sizes only the images that did not ask for a width,
% which here is nearly all the screenshots. To change one specific image, set
% `width=` on it in the .qmd, next to the picture it applies to.
\newcommand{\rbpfigscale}{0.7}

% \makebox keeps a less-than-full-width image centred. Without it the image
% sits hard against the left edge with a river of white space down the right,
% and `\centering` cannot help: an inline image comes wrapped in
% \pandocbounded rather than in a float, so there is no float to centre it in.
\usepackage{letltxmacro}
\LetLtxMacro{\rbporigincludegraphics}{\includegraphics}
\renewcommand{\includegraphics}[2][]{%
  \makebox[\linewidth]{%
    \rbporigincludegraphics[width=\rbpfigscale\linewidth,height=\maxdimen,keepaspectratio,#1]{#2}%
  }%
}

% --- The chapter overview box ----------------------------------------------
%
% overview-box.lua turns `::: {.overview}` into this environment, with the
% heading lifted out to become the title. The filter also carries a plain
% fallback definition for any PDF format that does not load this file, so this
% is the nice version rather than the only version.
%
% tcolorbox is loaded here rather than left to Quarto so the definition can sit
% in the preamble. Quarto emits
% `\@ifpackageloaded{tcolorbox}{}{\usepackage[skins,breakable]{tcolorbox}}`
% further down, which takes the empty branch once we have loaded it - and the
% options must match Quarto's exactly, or that becomes an option clash.
%
% The keys mirror Quarto's own callouts so the two sit together on a page,
% minus the colour: the five callout types each mean something specific and are
% navigated by colour, and the overview is furniture rather than a sixth
% meaning. Same reasoning as the HTML, in chapter-overview.scss.
\usepackage[skins,breakable]{tcolorbox}
\newtcolorbox{rbpoverview}[1]{%
  enhanced jigsaw,
  breakable,
  arc=.35mm,
  colback=white,
  opacityback=0,
  colframe=black!30,
  colbacktitle=black!6,
  opacitybacktitle=1,
  coltitle=black,
  left=2mm,
  leftrule=.15mm,
  rightrule=.15mm,
  toprule=.15mm,
  bottomrule=.15mm,
  titlerule=0mm,
  toptitle=1mm,
  bottomtitle=1mm,
  title={\textbf{#1}},
}

% --- Long code lines -------------------------------------------------------
% The body measure is narrower than it was, so an 80 character code line no
% longer fits. Break rather than run off the page.
\usepackage{fvextra}
\DefineVerbatimEnvironment{Highlighting}{Verbatim}{
  breaklines, breaknonspaceingroup, breakanywhere, commandchars=\\\{\}
}
